% Part: normal-modal-logic
% Chapter: tableaux
% Section: simple-S5

\documentclass[../../../include/open-logic-section]{subfiles}

\begin{document}

\olfileid{nml}{tab}{s5}

\olsection{\usetoken{P}{tableau}های ساده برای~\Log{S5}}

\Log{S5} نسبت به ردهٔ مدل‌های همگانی درست و کامل است؛ یعنی مدل‌هایی
که هر جهان در آن‌ها از هر جهان دسترسی‌پذیر است. در مدل‌های همگانی
رابطهٔ دسترسی‌پذیری اهمیتی ندارد: «جهان~$w$ای وجود دارد که
$\mSat{M}{!A}[w]$» صادق است اگر و تنها اگر چنین~$w$ای وجود داشته باشد
که از~$u$ دسترسی‌پذیر است. پس در~\Log{S5} می‌توانیم مدل‌ها را صرفاً
به‌صورت مجموعه‌ای از جهان‌ها و ارزش‌گذاری~$V$ تعریف کنیم. این نشان
می‌دهد که باید بتوانیم قواعد !!{tableau} را نیز ساده کنیم. در حالت کلی،
دنباله‌های اعداد صحیح مثبت را پیشوند می‌گیریم تا بتوانیم پیگیری کنیم
کدام پیشوندها جهان‌هایی را نام می‌برند که از جهان‌های دیگر
دسترسی‌پذیرند: $\sigma.n$ نام جهانی دسترسی‌پذیر از~$\sigma$ است. اما
در~\Log{S5} هر جهان از هر جهان دسترسی‌پذیر است، پس نیازی به این پیگیری
نیست. در عوض می‌توانیم از اعداد صحیح مثبت به‌عنوان پیشوند استفاده کنیم.
قواعد ساده‌شده در~\olref{tab:rules-S5} آمده‌اند.

\begin{table}
  \begin{center}
    \def\arraystretch{3}\def\fCenter{}
    \begin{tabular}{|c|c|}
    \hline
    \iftag{prvBox}{
      \AxiomC{\sFmla{\True}{\Box !A}[n]}
      \RightLabel{\TRule{\True}{\Box}}
      \UnaryInfC{\sFmla{\True}{!A}[m]}
      \DisplayProof
      &
      \AxiomC{\sFmla{\False}{\Box !A}[n]}
      \RightLabel{\TRule{\False}{\Box}}
      \UnaryInfC{\sFmla{\False}{!A}[m]}
      \DisplayProof\\
      $m$ استفاده‌شده است & $m$ تازه است
      \\[1ex]
      \hline}{}
    \iftag{prvDiamond}{
      \AxiomC{\sFmla{\True}{\Diamond !A}[n]}
      \RightLabel{\TRule{\True}{\Diamond}}
      \UnaryInfC{\sFmla{\True}{!A}[m]}
      \DisplayProof
      &
      \AxiomC{\sFmla{\False}{\Diamond !A}[n]}
      \RightLabel{\TRule{\False}{\Diamond}}
      \UnaryInfC{\sFmla{\False}{!A}[m]}
      \DisplayProof\\
      $m$ تازه است & $m$ استفاده‌شده است
      \\[1ex]
      \hline}{}
    \end{tabular}
  \end{center}
  \caption{قواعد ساده‌شده برای~\Log{S5}.}
  \ollabel{tab:rules-S5}
\end{table}

\begin{ex}
  تابلوی بستهٔ ساده‌شده‌ای ارائه می‌کنیم که نشان می‌دهد
  $\Log{S5} \Proves \Ax{5}$؛ یعنی~$\Diamond!A \lif \Box\Diamond!A$.
  \begin{oltableau}
    [\pFmla{\False}{\Diamond\formula{A} \lif \Box\Diamond \formula{A}}{1},
      just = \TAss
      [\pFmla{\True}{\Diamond \formula{A}}{1}, just = {\TRule{\False}{\lif}[1]}
        [\pFmla{\False}{\Box\Diamond \formula{A}}{1},
          just = {\TRule{\False}{\lif}[1]}
          [\pFmla{\False}{\Diamond \formula{A}}{2},
            just = {\TRule{\False}{\Box}[3]}
            [\pFmla{\True}{\formula{A}}{3},
              just = {\TRule{\True}{\Diamond}[2]}
                [\pFmla{\False}{\formula{A}}{3},
                  just = {\TRule{\False}{\Diamond}[4]}, close]
            ]
          ]
        ]
      ]
    ]
  \end{oltableau}
\end{ex}

\end{document}
